\section{A unitary non-Weil rank cutoff with every partner retained}
\label{sec:unitary-nonweil-point-cutoff}

\paragraph{Status and provenance.}
D-TD68, question Q2fx, is an accepted arbitrary-partner reduction.
The complete author proof/source reading and first exact supplement pass;
complete root and independent whole-proof reviews also pass. The
scalar-field summed-genus lemma in D-TD66, Section~1, is now accepted
at the exact reviewed mathematics; its immutable author freeze is retained.
The character-level and degree inputs are accepted results. The parent
suggested the all-projective-point comparison, and G34 supplied the
general scalar-field genus proof. The outer unitary centralizer argument
below combines a freshly read published correspondence with elementary
matrix counting and the accepted commutator average.

Put
\[
 D_0=10^{64},\qquad T_0=2^{1100000},\qquad B_0=10^{30},
 \quad C_0=2\cdot10^{30},\quad J_0=2C_0+1,
 \quad \sigma_0=10^{-29},\quad \varepsilon_q=q^{-B_0}.
 \tag{UPL.1}
\]
The unitary parameter is \(q=p^f\); the natural scalar field is
\(\mathbf F_Q\), where \(Q=q^2\). Natural geometric support means
codimension of the largest eigenspace over an algebraic closure.
Complex projective support allows every target eigenvalue.

\begin{theorem}[D-TD68]
Let an actual connected finite \(J\)-Galois sphere cover have faithful
irreducible complex coefficient module \(V\), of degree \(e\), with
generating product-one inertia operators satisfying
\[
                       \sum_i\operatorname{rank}(1-T_i)=2e+6.
 \tag{UPL.2}
\]
Suppose \(V=W\otimes U\) is an actual stable whole factorization.
The finite projective image on \(W\) contains and normalizes the
projective image of a finite perfect central product
\(E_1\cdots E_t\) acting faithfully and irreducibly on \(W\).
Its commuting quasisimple components are permuted transitively.
Each \(E_j\) is an actual central quotient of
\(\operatorname{SU}_d(q)\), with \(d\ge D_0\), and the faithful
irreducible component module \(W_j\) inflates to an ordinary character
of unitary level at least two. Write
\(W=\bigotimes_{j=1}^tW_j\), with common component degree \(a\).
If \(1\le t\le T_0\), this configuration is impossible.

The entire external partner \(U\), including its dimension, image and
scalar phases, is unrestricted. All prime powers \(q\ge2\), all
compatible component permutations and full field/diagonal normalizers
are covered. No projected-mark minimum, self-duality, character-field
equality or rationality of an intermediate base is assumed.
The named-cover premise does not include an arbitrary exceptional
central cover merely because its simple quotient is unitary.

Consequently, in exactly the primitive minimal naturally nonsplit
exact-imaginary-quadratic source class of D-TD54, every such non-Weil
unitary layer in a fixed complete D-TD12 recursion has \(d<D_0\).
Unitary level one is not excluded by this theorem.
\end{theorem}

\subsection*{1. The actual degree and level bounds}

Transitivity of the actual component action makes all component degrees
equal. Applying D-TD58 to this stable whole factor, with \(U\) still
present, excludes the simultaneous opposite degree inequalities. Thus
\[
                              a<q^{C_0d}.
 \tag{UPL.3}
\]
For an inflated irreducible \(\operatorname{SU}_d(q)\)-character of
level \(j\), the complete invoked GLT results, as recorded in
D-UNW-01, give
\[
 a\ge\frac{q^{j(d-j)}}{2(q+1)},\qquad
 j\ge d/2\Longrightarrow a>\frac{q^{d^2/4-2}}{q+1}.
 \tag{UPL.4}
\]
These are \cite[Definition~2, Lemma~7.1 and Theorem~7.3,
pp.~51--53]{UNW:GLT2020}. They and (UPL.3) imply
\(2\le j\le J_0<d/2\): the second inequality first excludes
\(j\ge d/2\), and then
\(a\ge q^{jd/2-3}\) bounds \(j<2C_0+6/d\).
Also, as proved from (UPL.4) in D-UNW-02,
\[
                         a\ge q^{2d-7}>q^d.
 \tag{UPL.5}
\]
Corollary~8.6(i) of GLT, including its proof, shows that a character
above this SU character restricts irreducibly, with unchanged degree.
After a determinant twist it is the unique true-level-\(j\) character
\(D_\alpha^\circ\) described in Section~8.3 of that paper.
The condition \(j^2+j\le d-1\) is amply satisfied. The complete
nonzero, irreducibility and exhaustion argument there works for even
as well as odd \(q\) in this range. We obtain an ordinary extension
\(\widetilde\chi\) to \(\operatorname{GU}_d(q)\).

Here and below GU denotes the isometry group of a nondegenerate
Hermitian form, as in GLT, not an additional arbitrary similitude
group. A unitary similitude can be rescaled to an isometry by
surjectivity of \(\mathbf F_{q^2}^*\to\mathbf F_q^*\),
\(c\mapsto c^{q+1}\); this does not change conjugation.
The norm and determinant assertions are also proved in
\cite[Lemma~10.1 and pp.~114--118]{OS:Taylor1992}, whose notation
\(U\), rather than \(GU\), denotes this same isometry group.

For \(g\in\operatorname{GU}_d(q)\) of geometric support \(s\),
GLT's exact (8.18), with (8.17), gives the following useful uniform
consequence:
\[
                 \frac{|\widetilde\chi(g)|}{a}
                         \le q^{-js}+10q^{-d/2}.
 \tag{UPL.6}
\]
Indeed, use the bounds \(|\operatorname{GU}_j(q)|\le2q^{j^2}\)
and \(\sqrt{16.52q^{j^2+j-1}}\le5q^{j^2}\). With
\(x=q^{j^2-d}\le q^{-d/2}\) and \(u=2x+10x^2\), the normalized
denominator in (8.18) is at least \(1-u\), and its numerator is at
most \(q^{-js}+u\). Since \(x<1/100\), one has
\(u\le3x\le3/100\), and
\(2u/(1-u)\le600x/97<7x\). This proves the weaker displayed
constant ten. For central \(g\) the displayed bound is automatic.
All estimates are independent of the chosen determinant twist.

\subsection*{2. All projective points and the summed genus}

Let \(\Lambda\) be all \(\mathbf F_Q\)-projective lines of the
natural space, with
\[
                         N=|\Lambda|=\frac{Q^d-1}{Q-1}.
 \tag{UPL.7}
\]
For a projective linear operator \(g\), write \(v(g)\) for the
codimension of its largest rational \(\mathbf F_Q\)-eigenspace,
using dimension zero if there is none. Elementary eigenspace counting
gives
\[
                 f_\Lambda(g)\le Q^{-v(g)}+2Q^{1-d/2}.
 \tag{UPL.8}
\]
For completeness, if its largest rational eigenspace has dimension
\(k\ge d/2\), the other eigenspaces have total dimension at most
\(d-k\); their line counts sum to at most the count in dimension
\(d-k\). Thus the fixed fraction is at most
\((Q^k+Q^{d-k}-2)/(Q^d-1)\le Q^{k-d}+2Q^{-d/2}\).
If \(k<d/2\), at most \(Q-1\) rational eigenvalues occur, giving
the weaker upper bound \(2Q^{1-d/2}\).

Conversely, if geometric support \(s<d/2\), its largest eigenvalue
is rational: any distinct Galois conjugate would supply another
eigenspace of the same dimension, contradicting the natural dimension.
Its eigenspace therefore contributes at least
\((Q^{d-s}-1)/(Q^d-1)\ge Q^{-s}-2Q^{-d}\) to the fixed fraction.
If \(s\ge d/2\), \(q^{-js}\le q^{-d}\) because \(j\ge2\).
Combining both cases with (UPL.6) proves
\[
           \frac{|\widetilde\chi(g)|}{a}
             \le f_\Lambda(g)+12q^{-d/2}
             \le f_\Lambda(g)+\varepsilon_q.
 \tag{UPL.9}
\]
The level-two hypothesis is essential here; it is not replaced by a
level-one assertion.

We now invoke exactly the general scalar-field lemma proved in
D-TD66, Section~1. It applies to any finite
\(G\subset\operatorname{PGL}_d(Q)\) with a normal nonabelian simple
subgroup \(S\) of order greater than 60 and irreducible natural
linear preimage, and any nonempty finite \(G\)-set satisfying
\(f(g)\le Q^{-v(g)}+1/10000\) for \(g\ne1\).
The set need not be transitive. For the point-stabilizer quotient
curves of every actual connected \(G\)-cover of any base curve,
the lemma gives
\[
                         \sum_\beta(g(Y_\beta)-1)
                                  \ge\frac{|\Lambda|}{10000}.
 \tag{UPL.10}
\]
The proof uses the actual connected regular group-cover, weak
scalar-choice Scott inequalities and the summed Riemann--Hurwitz
formula. Over an elliptic base it uses forced nontrivial branching;
over a higher-genus base the claim follows directly from that formula.
It does not require an intermediate base to be rational.

Take now \(\operatorname{PSU}_d(q)\le G\le\operatorname{PGU}_d(q)\).
Its normal subgroup \(S=\operatorname{PSU}_d(q)\) is simple of order
greater than 60 in our range. Its natural linear preimage is
irreducible, even absolutely: in an orthonormal basis SU contains the
determinant-one diagonal group of norm-one scalars. Its distinct
coordinate characters, of order prime to the characteristic, make
every invariant subspace a sum of coordinate lines. The even
permutation matrices also lie in SU and act transitively on those
lines. Hence only zero and the full space are invariant.
Since \(Q=q^2\ge4\) and \(2Q^{1-d/2}<1/10000\), (UPL.8)
checks the remaining hypothesis. Thus (UPL.10) holds for the all-line
set, with both unitary point types present. No transitivity assertion
on that set is used.

\subsection*{3. The genuine full unitary normalizer}

Write \(A=\operatorname{Aut}(S)\) and
\(B=\operatorname{PGU}_d(q)\). Steinberg's Theorem~3.2 and the
complete twisted-type proof in Section~7, together with (3.3)--(3.6),
show that
\[
 B\triangleleft A,\qquad c=[A:B]\le2f,
 \qquad \alpha(X)=M X^\sigma M^{-1},
 \quad M\in\operatorname{GU}_d(q),\quad
 \sigma\in\operatorname{Gal}(\mathbf F_{q^2}/\mathbf F_p).
 \tag{UPL.11}
\]
See \cite[pp.~607--608 and 613--614]{PFL:Steinberg1960}; the
supplied normal-type proof has already been read in D-TD65.
The inner-diagonal group is PGU: its index over PSU is
\(\gcd(d,q+1)\), the diagonal index in Steinberg's (3.4).
For unitary groups the graph operation is already the field map
\(X\mapsto X^q\), because in a standard Hermitian basis
\(X^q=X^{-T}\). Thus a field part of order two is exactly
transpose-inverse, even after the arbitrary diagonal adjustment
\(M\). There is no missing additional graph coset in (UPL.11).

These natural automorphisms also give the correct actual complex
normalizer interface for every allowed central quotient. A local
operator inducing \(\alpha\) on \(S\) is projectively an intertwiner
of the two inflated SU representations. Its scalar discrepancy is
a linear character of the perfect group SU, hence is trivial.
It is therefore an honest intertwiner after lifting \(\alpha\) by
(UPL.11). This argument also proves compatibility with the actual
central kernel; it does not assume that every automorphism of \(S\)
is compatible with a given character. When \(\alpha\in B\), the
local operator differs by a scalar from \(\widetilde\chi(M)\).
When an operator induces the identity on all projective components,
perfectness kills all its central discrepancies and Schur's lemma
makes it scalar. Hence the actual projective image embeds in the
component-permuting automorphism group of \(S^t\).

All trace calculations can be made unitary. A finite projective
matrix group has a finite determinant-one lift; averaging gives an
invariant positive Hermitian form. Local component intertwiners and
cycle products may then be chosen unitary. Multiplication by a
unit complex scalar leaves every absolute trace ratio unchanged.

\subsection*{4. Every outer unitary trace is exponentially small}

We first bound the projective centralizer of any \(\alpha\notin B\).
Let \(l\ge2\) be the order of its field part \(\sigma\) in
(UPL.11). A projectively fixed lift \(X\in\operatorname{SU}_d(q)\)
satisfies \(M X^\sigma M^{-1}=zX\), where \(z\) is a central
norm-one scalar. There are at most \(q+1\) possible scalars.

If \(l\ge3\), the equation for each scalar is a semilinear fixed
equation on the \(d^2\)-dimensional matrix space over
\(\mathbf F_{q^2}\). Its number of solutions is at most
\(q^{2d^2/l}\). Here is a direct proof that does not assume the
semilinear map itself has order \(l\). Vectors fixed by it and
independent over the fixed scalar field are independent over the
larger field. Otherwise choose a shortest dependence with one
coefficient one, apply the map and subtract. Minimality makes every
coefficient fixed by \(\sigma\), a contradiction. The fixed-space
dimension over the fixed scalar field is therefore at most \(d^2\).

If \(l=2\), then \(\sigma:x\mapsto x^q\), and the ordinary fixed
subgroup in GU is the centralizer of \(M\tau\), where
\(\tau(X)=X^{-T}\). Its intersection with SU bounds the relevant
ordinary fibers. Gow--Vinroot's Theorem~2.1 gives a corresponding
element \(D\tau\) in \(\operatorname{GL}_d(q)\tau\) with an
isomorphic full centralizer; this theorem holds for all \(q\) and
arbitrary coset-element orders. We use its complete supplied proof,
Lemmas~2.1--2.4 and Theorem~2.1 in
\cite[author pp.~3--6]{UPL:GowVinroot2008}.
Its Lang--Steinberg input is an invoked published theorem, not a
newly supplied proof here.

The latter centralizer is the isometry group of a nonsingular
bilinear matrix, with no symmetry hypothesis. Its order is at most
\(q^{d(d+1)/2}\): after \(j-1\) independent columns have been
chosen, the first \(j-1\) pairing equations are independent linear
conditions on column \(j\), leaving at most \(q^{d-j+1}\) choices.
Multiply over \(j\), ignoring additional equations. Each nonempty
projective scalar-discrepancy fiber is a coset of the ordinary
fixed subgroup, so the extra factor \(q+1\) suffices.

Finally, the complete order calculation in
\cite[pp.~117--118]{OS:Taylor1992} gives
\(|\operatorname{SU}_d(q)|\ge(2/3)q^{d^2-1}\): in its product
formula, the factors with odd index exceed one, while the product
of the even-index factors is at least
\(1-\sum_{j\ge1}q^{-2j}\ge2/3\).
The central quotient cancels when taking the fixed proportion.
Both cases therefore give
\[
 \frac{|C_S(\alpha)|}{|S|}
   \le3q^{-d^2/3+2};\qquad
 l=2:\quad \frac{|C_S(\alpha)|}{|S|}
   \le3q^{-d^2/2+d/2+2}.
 \tag{UPL.12}
\]
The second bound implies the first for \(d\ge3\).

Put \(R=\lfloor d/32\rfloor\). At most
\(q^{4dR+2}\) elements of \(S\) have geometric support below
\(R\). Indeed their dominant eigenvalue is rational, and a scalar
shift of a lift has rank at most \(R\); matrices of that rank are
overcounted by products of a \(d\)-by-\(R\) and an
\(R\)-by-\(d\) matrix over \(\mathbf F_{q^2}\), with at most
\(q^2\) scalar choices. Every nonempty fiber of
\(x\mapsto x^{-1}\alpha(x)\) has size \(|C_S(\alpha)|\).
Thus the fraction of its low-support values, counted in its domain,
is at most
\[
        3q^{4dR-d^2/3+4}
             \le3q^{-5d^2/24+4}\le q^{-d^2/6}.
 \tag{UPL.13}
\]
The last inequality already holds for \(d\ge12\).

For a unitary local intertwiner \(T\), Schur averaging over the
represented perfect core gives
\[
 \left|\frac{\operatorname{tr}T}{a}\right|^2
 \le\frac1{|S|}\sum_{x\in S}
       \frac{|\chi(x^{-1}\alpha(x))|}{a}
 \le q^{-d^2/6}+q^{-\sigma_0R}.
 \tag{UPL.14}
\]
Absolute character values are independent of the chosen central
lift. The high-support bound is the accepted all-element LT estimate
\(a^{-\sigma_0s/d}\), strengthened using (UPL.5).
Since \(R\ge d/64\), both exponents on the right are at least
\(2B_0+1\). Therefore
\[
                         |\operatorname{tr}T|/a\le\varepsilon_q
                              \qquad(\alpha\notin B).
 \tag{UPL.15}
\]
No restriction on the order of the outer operator or its original
inertia lift was made.

\subsection*{5. The actual frame-and-point cover}

Use the genuine induced \(A\)-set
\(A\times_B\Lambda\). It has \(c\) frames of \(N\) points.
An element outside \(B\) fixes no frame, since \(B\triangleleft A\)
and the action on \(A/B\) is regular. For an element in \(B\),
the fixed fraction is the same in every frame: its conjugates by
the field automorphisms in (UPL.11) have identical rational
eigenspace dimensions. The all-line action of \(B\) is faithful;
the induced \(A\)-action is consequently faithful as well.

Transport these sets through the actual component identifications.
The original group has an actual finite permutation action on
\[
               \Omega=\coprod_{j=1}^t(A\times_B\Lambda),
                             \qquad |\Omega|=ctN.
 \tag{UPL.16}
\]
There is no assertion of a separate original quotient to a permuted
individual component. For each orbit on component frames, take its
actual stabilizer and corresponding intermediate connected base
\(X\). Its local linear image lies between PSU and PGU and
contains PSU, because the selected actual component core is present.
The point covers above that frame, including all their orbits, have
summed genus surplus at least \(N/10000\) by (UPL.10).
They are exactly the point-stabilizer quotient curves for the
corresponding part of \(\Omega\) over the original sphere.
There is at least one frame orbit. Summing their actual
Riemann--Hurwitz equalities gives, for the original inertia images
\(y_i\),
\[
 \sum_i\frac{\operatorname{ind}_\Omega(y_i)}{|\Omega|}
          \ge2+\frac{2}{10000ct}
          \ge2+\frac1{10000ft}.
 \tag{UPL.17}
\]
This step uses summed genera, not a false transitivity claim for
the two unitary point types.

\subsection*{6. Whole tensor traces and original-period Fourier}

For a component-permuting operator on \(W\), the usual tensor-cycle
trace identity gives normalized trace equal in absolute value to
\(a^{-I}\) times the product of the normalized local cycle traces,
where \(I\) is the permutation index on components.
If there are no fixed components, this is at most \(a^{-1}<
\varepsilon_q\), and no point of \(\Omega\) is fixed.
If a fixed component has an outer operator, (UPL.15) gives the
same bound for the whole product.

Otherwise let \(b\ge1\) be the number of fixed components, all
inner-diagonal. Their product of normalized traces is at most their
average. Equation (UPL.9), followed by
\(a^{-I}\le b/t\), proves
\[
                  \frac{|\operatorname{tr}A|}{\dim W}
                           \le f_\Omega(y)+\varepsilon_q.
 \tag{UPL.18}
\]
For the elementary factor inequality, if the component permutation
is trivial then \(b=t,I=0\). Otherwise it moves \(k=t-b\ge2\)
components, \(I\ge k/2\), and \(a\ge4\) gives
\(a^I\ge2^k\ge1+k\ge t/b\).
All remaining cycle traces have absolute normalized value at most
one. If the natural action is the identity, the operator on \(W\)
is projectively scalar by the faithful embedding and the faithful
set action, and (UPL.18) holds with error zero.

The accepted full-normalizer support theorem FT20--21 gives
\(\nu(A)/\dim W\ge1/80\) for every nonidentity projective operator
on this entire unitary layer; a moved component also has the
stronger quarter bound. There is no alternating component here, by
the full repetition list in
\cite[author Section~2.2.1, p.~13]{HES:LiXia2022}.
Tensor spectral convolution cannot increase a normalized largest
eigenspace, so every original entry with \(y_i\ne1\) costs at
least \(e/80\). Since \(e\ge a>q^d>480\), (UPL.2) implies
that there are at most160 such entries.

Retain now each full original operator order \(L_i\).
Write \(T_i=A_i\otimes B_i\) and normalize reciprocal unitary
phases so that both factors have \(L_i\)-th power one.
The natural image has order dividing \(L_i\), and
(UPL.18) applies to every power of \(A_i\), including all
projectively identity powers. The absolute trace of the normalized
partner factor is at most one. Fourier averaging at target one,
and permutation Cauchy--Frobenius averaging over the same full
period, therefore give
\[
 \frac{\operatorname{rank}(1-T_i)}e
       \ge\frac{\operatorname{ind}_\Omega(y_i)}{|\Omega|}
               -\varepsilon_q\quad(y_i\ne1),
 \qquad
 \frac{\operatorname{rank}(1-T_i)}e\ge0\quad(y_i=1).
 \tag{UPL.19}
\]
No scalar period has been replaced by a shorter natural period.
Summing (UPL.19) and (UPL.17) would imply
\(2+6/e\ge2+1/(10000ft)-160\varepsilon_q\).
But \(6/e<6\varepsilon_q\), while
\[
 166f\varepsilon_q\le166q^{1-B_0}
     \le166\,2^{1-B_0}<\frac1{10000\,2^{1100000}}
     \le\frac1{10000t}.
 \tag{UPL.20}
\]
Here \(f\le q\), and the last strict comparison is an integer
power-of-two inequality. This contradiction proves the standalone
statement.

For the stated main-source corollary, use precisely the accepted
D-TD54 primitive minimal naturally nonsplit class. Its complete
D-TD12 system has whole support denominator
\(M_0=\max(80,N_*-2)\), and D-TD16 gives \(t\le M_0/2\).
The exact estimate \(N_*<2^{1049312}<2^{1100000}\), already
checked in D-TD65, supplies the needed bound on \(t\).
All other standalone hypotheses are the actual named-cover and
level premises of the selected layer. No new transfer to a broader
quasiprimitive class, no level-one cutoff, and no assertion about
source existence is made.
